%% sections/01-introduction.tex

\section{Introduction}
\label{sec:introduction}

Tabletop role-playing games (TRPGs) are a site of rich, nuanced player behavior, yet
that behavior resists systematic description. At most tables, sessions are unrecorded.
Player choices are attributed after the fact to player personality, mood, or the
chemistry of a particular group, rather than to identifiable features of the character
and campaign design. When a player does something surprising and resonant --- names a
grief their character has been carrying, performs a small act of care that changes the
scene's tone, or produces a speech at the finale that feels like an inevitable
conclusion --- there is no formal mechanism for capturing what enabled it, let alone
for replicating it in future designs.

The absence of such a mechanism is a design problem, not merely an ethnographic one.
TRPG character-sheet design has matured considerably in terms of mechanical specificity
(stat blocks, action economies, spell lists) and narrative scaffolding (backstory
prompts, bond/flaw fields, ideals). What it has not developed is a vocabulary for
\textit{behavioral design affordances}: specific character-sheet features that reliably
enable specific classes of behavior during play, independent of who is playing. Without
such a vocabulary, designers cannot learn from surprising sessions or build on them.
Every table starts from scratch.

\subsection{The Contribution}

This paper describes a pipeline --- developed and validated in the Marathon D\&D
Workshop --- for converting session-level behavioral surprises into ratified design
patterns. The pipeline is:

\begin{enumerate}
  \item \textbf{Innovation logging.} During and after each session, a session-gate
    reviewer logs behavioral surprises: moments when a player character's action or
    speech exceeds what the existing rubric anticipates. These are \textit{innovation
    instances}.
  \item \textbf{Threshold accumulation.} When 3 or more innovation instances cluster
    thematically across 2 or more sessions, the cluster crosses the promotion threshold.
    The threshold is conservative by design: it requires cross-session evidence (the
    pattern is not session-specific) and cross-instance evidence (the pattern is not a
    single-event anomaly).
  \item \textbf{Player-style proposal.} A player-style entry is drafted. The entry
    includes a definition in behavioral terms (not player-personality terms), observable
    session-log indicators, design features that enable the style, and a cross-campaign
    validation status.
  \item \textbf{Ratification.} The style is ratified when a second review confirms that
    the promotion-threshold instances are substantively distinct and that the definition
    correctly names the observed pattern.
  \item \textbf{Cross-campaign validation.} After ratification, the style is monitored
    across subsequent campaigns to confirm it fires outside the archetype that produced
    it.
\end{enumerate}

Applied across 49 sessions and 7 complete campaigns, the pipeline promoted 7 player
styles and validated them against 12 research hypotheses (RV1--RV12). The hypotheses
range from basic cross-campaign recurrence (RV1: sheet-deep-reader fires outside C1)
to archetype-independence (RV9: styles fire without any designed personal stakes) to
trigger-type shift (RV11: trigger type shifts from situational to relationship-driven
across a campaign's arc).

\subsection{Why AI Simulation Is a Useful Laboratory}

Human TRPG sessions produce behavioral patterns that are difficult to observe
systematically for several reasons. Play is real-time and typically unrecorded.
Player choices blend design input (what the character sheet affords) with player
personality (the individual at the table) and social dynamics (the chemistry of the
group). Separating the contribution of character-sheet design from those other factors
requires either controlled experiments that eliminate ecological validity or longitudinal
naturalistic observation that is difficult to conduct at scale.

In the Marathon corpus, AI agents play player characters strictly according to their
character-sheet Playstyle Heuristics. There is no player personality to confound with
sheet design. If a behavioral pattern fires repeatedly, it fires from the sheet. This
does not mean that human players would not produce the same patterns --- we expect
they would, and in greater variety. What it means is that the AI corpus isolates the
design contribution: any pattern that fires in the AI corpus is demonstrably enabled
by character-sheet design rather than by player personality.

This is the key methodological advantage. It comes with corresponding limitations.
AI-simulated play adheres to heuristics more consistently than human play, which
introduces a regularity that human sessions would not always match. The AI cannot
improvise beyond its heuristics in the way that a skilled human player can. And the
corpus represents a single AI system, making generalization to other systems an open
question. These limitations are discussed in Section~\ref{sec:discussion}.

\subsection{The Player-Style vs.\ Player-Type Distinction}

Player typologies in the literature describe \textit{who plays games and why}. Bartle's
four types (Achiever, Explorer, Socializer, Killer) and Yee's motivation dimensions
(Achievement, Social, Immersion and their sub-components) characterize players as
people: their motivations, preferences, and psychological needs~\cite{bartle1996hearts,
yee2006motivations}. These are valuable for understanding audience composition and
long-term engagement, but they do not tell a designer what to put in a character sheet
to enable a specific class of behavior.

Player styles, as defined here, describe \textit{what a character sheet enables a PC
to do in a session}. Sheet-deep-reader is not a player type (``the sort of person who
reads their sheet carefully''); it is a behavioral design affordance (``a PC sheet
with a detailed grief-paragraph and local-origin notes enables the PC to surface
sheet-content as current-scene inputs without DM prompting''). The design implication
is direct: a designer who wants sheet-deep-reader behavior writes a denser, more
specific character sheet. A designer who wants Achiever behavior cannot design that
in the same way.

\subsection{Roadmap}

Section~\ref{sec:related} reviews related work in player typologies, emergent narrative,
AI player behavior, and design patterns in games. Section~\ref{sec:methodology} describes
the promotion pipeline in full, including the seven promoted styles with definitions,
instance counts, and design triggers. Section~\ref{sec:evaluation} presents the
cross-campaign validation results for RV1--RV12, including per-campaign instance tables
and trigger-type breakdown. Section~\ref{sec:discussion} addresses the significance
of the findings and the limitations of the corpus. Section~\ref{sec:conclusion}
concludes and identifies directions for future work.
